\begin{pastexam}[label={px:gk3-115-ds-01}]{高考三級資訊處理 115 年 第一題}
（申論；25 分）

某系統 A 使用雜湊表（hash table）儲存不同的正整數鍵值，亦即不允許重複鍵值，雜湊表有 11 個儲存格，索引從 0
開始，雜湊函數（hash function）為$h^A(k) = k \bmod 11$，其用平方探查法（quadratic probing）處理碰撞
（collision）問題，探查序列為 $h_i^A(k) = (h^A(k) + i^2) \bmod 11\ (i = 0, 1, 2, \ldots)$，刪除資料時，被刪除
資料的位置標記為特殊符號 DELETED。請回答下列問題：
\begin{itemize}[leftmargin=2.2em]
\item[㈠] 給定一個空的雜湊表，依序插入下列鍵值 22、1、13、24、35、46、7、18，請畫出所有鍵值插入完成後的雜湊
表狀態，並列出插入鍵值 46 時的完整探查過程。（10 分）
\item[㈡] 承上題，依序刪除鍵值 24、13，畫出刪除後的雜湊表狀態。並說明為什麼刪除鍵值時需用 DELETED 標記，而不
能將該鍵值所在的儲存格恢復成「從未存放過鍵值」的空狀態。（5 分）
\item[㈢] 承上題，執行插入鍵值 12，請列出插入時的探查過程、操作停止的理由，並寫出 12 最後插入那一個儲存格。
插入時，DELETED 標記視為可放入新鍵值的儲存格。請注意鍵值不能重覆。（5 分）
\item[㈣] 相較於系統 A，考慮另一個採用平方探查法之系統，系統 B 的表格大小為 8，索引亦從 0 開始，雜湊函數
$h^B(k)$ 及探查序列 $h_i^B(k)$ 分別定義為 $h^B(k) = k \bmod 8$、
$h_i^B(k) = (h^B(k) + i + 2 \cdot i^2) \bmod 8\ (i = 0, 1, 2, \ldots)$。在非均勻雜湊（non-uniform hashing）的情況
下，也就是許多鍵值可能被分配到相同或少數幾個初始雜湊位置時，那一個系統的雜湊表儲存格利用率可能較高？
並說明理由。（5 分）
\end{itemize}
\end{pastexam}

\begin{pastsoln}
答案：㈠ 表為 $[22, 1, 13, 24, -, 18, 35, 46, 7, -, -]$（索引 $0 \sim 10$）；$46$ 依序探查 $2, 3, 6, 0, 7$，放在
$7$。 ㈡ $[22, 1, \text{D}, \text{D}, -, 18, 35, 46, 7, -, -]$（D 表示 DELETED）。 ㈢ 探查 $1, 2, 5, 10$，在空格
$10$ 停止，$12$ 放進儲存格 $2$。 ㈣ 系統 B。

\textbf{㈠} 每個鍵先算 $h = k \bmod 11$，再依序試 $h, h + 1, h + 4, h + 9, h + 16, \ldots \pmod{11}$，找到空格就
放（資結書第~12.3~節）。
\begin{center}
\small
\begin{tabular}{cccl}
\toprule
鍵 $k$ & $k \bmod 11$ & 探查的儲存格（$i = 0, 1, \ldots$） & 放在 \\
\midrule
22 & 0 & 0 & 0 \\
1 & 1 & 1 & 1 \\
13 & 2 & 2 & 2 \\
24 & 2 & 2（有 13）, 3 & 3 \\
35 & 2 & 2, 3, 6 & 6 \\
46 & 2 & 2, 3, 6, 0, 7 & 7 \\
7 & 7 & 7（有 46）, 8 & 8 \\
18 & 7 & 7, 8, 0, 5 & 5 \\
\bottomrule
\end{tabular}
\end{center}
\textbf{插入 46 的完整探查過程}：$h^A(46) = 46 \bmod 11 = 2$。
\begin{itemize}
\item $i = 0$：$(2 + 0) \bmod 11 = 2$，已有 $13$，碰撞。
\item $i = 1$：$(2 + 1) \bmod 11 = 3$，已有 $24$，碰撞。
\item $i = 2$：$(2 + 4) \bmod 11 = 6$，已有 $35$，碰撞。
\item $i = 3$：$(2 + 9) \bmod 11 = 0$，已有 $22$，碰撞。
\item $i = 4$：$(2 + 16) \bmod 11 = 7$，空的，放入 $46$。
\end{itemize}
（$18$ 的探查：$7 \to 8 \to (7 + 4) \bmod 11 = 0 \to (7 + 9) \bmod 11 = 5$。）全部插入後：
\begin{center}
\small
\begin{tabular}{l*{11}{|c}|}
\cline{2-12}
索引 & 0 & 1 & 2 & 3 & 4 & 5 & 6 & 7 & 8 & 9 & 10 \\ \cline{2-12}
鍵 & 22 & 1 & 13 & 24 & & 18 & 35 & 46 & 7 & & \\ \cline{2-12}
\end{tabular}
\end{center}

\textbf{㈡} 刪除 $24$：從 $h = 2$ 開始找，$2$ 是 $13$，$3$ 是 $24$，找到，標成 DELETED。 刪除 $13$：$2$ 就是
$13$，標成 DELETED。
\begin{center}
\small
\begin{tabular}{l*{11}{|c}|}
\cline{2-12}
索引 & 0 & 1 & 2 & 3 & 4 & 5 & 6 & 7 & 8 & 9 & 10 \\ \cline{2-12}
鍵 & 22 & 1 & D & D & & 18 & 35 & 46 & 7 & & \\ \cline{2-12}
\end{tabular}
\end{center}
\textbf{為什麼要用 DELETED}（資結書第~12.3.4~節）：開放定址的搜尋是沿探查序列一格一格找，遇到「從未存放過」
的空格就斷定鍵不存在而停止。 一個鍵當初被放到後面的位置，是因為前面的格子在那時都有人。 如果把刪除的格子
改回空格，探查序列會在那裡提早斷掉。 例如把儲存格 $2$、$3$ 改回空的，之後搜尋 $35$ 或 $46$（$h = 2$）時，
第一步就碰到空格 $2$，會誤判「不在表中」，但它們其實在 $6$、$7$。 DELETED 的意思是「這裡現在沒有鍵，但曾經
有過，搜尋要繼續往下找」；插入時則可以把新鍵放在這種格子。

\textbf{㈢ 插入 12}：$h^A(12) = 12 \bmod 11 = 1$。 因為不允許重複，插入前要確定 $12$ 不在表中，所以要一路探查到
「從未存放過鍵值」的空格才能下結論；途中記下第一個可用（DELETED）的格子。
\begin{itemize}
\item $i = 0$：儲存格 $1$，是 $1$，不是 $12$，繼續。
\item $i = 1$：儲存格 $(1 + 1) \bmod 11 = 2$，DELETED，記下「第一個可放的位置是 $2$」，但 $12$ 可能在後面，繼續。
\item $i = 2$：儲存格 $(1 + 4) \bmod 11 = 5$，是 $18$，繼續。
\item $i = 3$：儲存格 $(1 + 9) \bmod 11 = 10$，從未存放過鍵值的空格。
\end{itemize}
\textbf{停止的理由}：碰到空格 $10$，表示 $12$ 若存在一定在這之前出現過，所以 $12$ 不在表中，不會重複，探查
結束。 \textbf{放在}記下的第一個可用位置：儲存格 $2$（DELETED）。 結果：
$[22, 1, 12, \text{D}, -, 18, 35, 46, 7, -, -]$。 （若在儲存格 $2$ 就直接放入，萬一 $12$ 早已存在後面的格子，就會
重複。）

\textbf{㈣ 系統 B 的利用率可能較高。} 關鍵是同一個初始位置 $h$ 出發的探查序列能走到幾個不同的格子。
\begin{itemize}
\item \textbf{系統 A}：位移是 $i^2 \bmod 11$，$i = 0, 1, \ldots, 10$ 得到 $0, 1, 4, 9, 5, 3, 3, 5, 9, 4, 1$，只有
$\{0, 1, 3, 4, 5, 9\}$ 共 $6$ 種（奇質數 $m$ 的平方數模 $m$ 只有 $(m + 1)/2$ 種；資結書定理~12.3 說前
$(m+1)/2$ 個位置兩兩不同，之後只是重複）。 所以同一個初始位置的鍵最多只能用到 $11$ 格中的 $6$ 格。 例如初始
位置 $2$ 的鍵只能放在 $2, 3, 6, 0, 7, 5$，永遠用不到 $1, 4, 8, 9, 10$；許多鍵擠在同一個初始位置時，第 $7$ 個就
找不到位置，即使表中還有 $5$ 個空格。
\item \textbf{系統 B}：位移是 $(i + 2i^2) \bmod 8$，$i = 0, 1, \ldots, 7$ 得到 $0, 3, 2, 5, 4, 7, 6, 1$，恰好是
$0 \sim 7$ 每個各一次。 所以從任何初始位置出發，前 $8$ 次探查會走遍全部 $8$ 個格子，只要表沒滿就一定插得進去。
\end{itemize}
在非均勻雜湊下，很多鍵的初始位置相同，它們的命運完全由「一條探查序列能用幾格」決定：系統 A 最多用到表的
$6/11 \approx 55\%$，系統 B 可以用到 $100\%$。 所以系統 B 的儲存格利用率較高。 （$i + 2i^2$ 能走遍是因為
$i + 2i^2 \equiv i(2i + 1)$，對 $m = 2^3$ 而言 $i \mapsto i + 2i^2$ 是 $\{0, \ldots, 7\}$ 上的一對一函數：若
$i + 2i^2 \equiv j + 2j^2 \pmod 8$，則 $(i - j)(1 + 2(i + j)) \equiv 0$，而 $1 + 2(i + j)$ 是奇數、在模 $8$ 下
可逆，所以 $i \equiv j$。）
\end{pastsoln}
